\section{Metodologia}

A metodologia adotada para o desenvolvimento da plataforma Wumpus Verse baseou-se na separação entre frontend, backend e lógica dos agentes, seguindo uma arquitetura cliente-servidor de três camadas. Esta seção descreve os requisitos estabelecidos para o sistema, a arquitetura e tecnologias empregadas, os agentes implementados, os mecanismos de persistência de dados e a interface desenvolvida para a plataforma.

\subsection{Requisitos, Arquitetura e Tecnologias} \label{sec:req-arq-tec}

Para orientar o desenvolvimento da plataforma Wumpus Verse, foram definidos requisitos funcionais e não funcionais (Tabela~\ref{tab:requisitos}) que garantem flexibilidade, personalização e boa experiência ao usuário, além de uma arquitetura e tecnologias que agilizem e tornem o processo eficiente.

\begin{table}[ht]
    \centering
    \small
    \begin{tabular}{|c|c|p{6.5cm}|}
        \hline
        \textbf{Código} & \textbf{Tipo} & \textbf{Descrição} \\
        \hline
        RF01 & Funcional & Criação de grids de qualquer formato (não apenas configurações quadradas). \\
        RF02 & Funcional & Definição de parâmetros, objetivos e funções de avaliação para agentes. \\
        RF03 & Funcional & Execução de agentes personalizados em ambientes salvos. \\
        RF04 & Funcional & Salvamento de resultados e histórico de cada simulação. \\
        RF05 & Funcional & Geração de relatórios com pontuação, indicadores e caminho percorrido. \\
        RF06 & Funcional & Exportação de dados em JSON (ambientes, agentes e execuções). \\
        \hline
        RNF01 & Não Funcional & Multiplataforma: funcionamento em navegadores modernos (Chrome, Firefox, Edge, Safari). \\
        RNF02 & Não Funcional & Responsividade parcial para dispositivos móveis (funcionalidades reduzidas). \\
        RNF03 & Não Funcional & Disponibilidade 24/7 durante o período de avaliação. \\
        RNF04 & Não Funcional & Tempo de resposta menos de 5 segundos por requisição. \\
        \hline
    \end{tabular}
    \caption{Requisitos funcionais e não funcionais da plataforma Wumpus Verse.}
    \label{tab:requisitos}
\end{table}

Com base nesses requisitos, a arquitetura da plataforma foi projetada seguindo o modelo cliente-servidor de três camadas, conforme ilustrado na Figura~\ref{fig:diagrama-sistema}. A escolha das tecnologias para cada camada considerou critérios de produtividade, integração e maturidade do ecossistema. A camada de frontend (I), desenvolvida com a biblioteca React, é responsável pela interface gráfica e pelo consumo da API REST. O React foi selecionado por oferecer componentização, Virtual DOM e um ecossistema maduro, fatores que agilizam o desenvolvimento de interfaces interativas. A camada de backend (II), implementada em Python com o framework FastAPI, gerencia as requisições, coordena as simulações e controla o acesso ao banco de dados. A combinação de Python com FastAPI foi escolhida por proporcionar agilidade no desenvolvimento e integração direta com os agentes (que também são implementados em Python) eliminando a sobrecarga de comunicação entre linguagens e permitindo o reuso de código entre a API e os agentes. Por fim, a camada de persistência (III) é composta por dois sistemas distintos: o Firebase Authentication, responsável pelo gerenciamento de usuários, e o PostgreSQL, que armazena os dados específicos da aplicação. Essa separação de responsabilidades, aliada à escolha tecnológica adequada, permite um desenvolvimento modular e facilita a manutenção e evolução do sistema.

\begin{figure}[ht]
    \centering
    \includegraphics[width=0.6\textwidth]{Diagrama Arquitetura.png}
    \caption{Diagrama da arquitetura do sistema.}
    \label{fig:diagrama-sistema}
\end{figure}

\subsection{Agentes}

A plataforma Wumpus Verse disponibiliza três tipos de agentes com diferentes estratégias de decisão, todos projetados para maximizar a pontuação final, seja coletando o ouro, eliminando o Wumpus ou retornando à saída. A implementação desses agentes permite que os usuários comparem diferentes abordagens para o mesmo problema, facilitando a compreensão das vantagens e limitações de cada estratégia.

O \textbf{Agente Reativo Aleatório} implementa uma abordagem simples baseada em regras. A cada passo, o agente recebe as percepções da sala atual. Se detecta ouro, o agente o coleta imediatamente. Se detecta fedor, existe uma chance de 10\% de o agente atirar ou mover-se em uma direção aleatória. Este agente ignora completamente a percepção de brisa, o que o torna vulnerável a cair em poços, mas serve como um baseline para comparação com estratégias mais sofisticadas.

O \textbf{Agente Baseado em Modelo} mantém um estado interno, ou crença, que é atualizado continuamente com base nas percepções recebidas. Salas consideradas seguras são exploradas, enquanto salas com indícios de perigo são evitadas. Este agente utiliza o algoritmo A* para planejar rotas ótimas entre sua posição atual e os destinos desejados. O agente pode ser configurado com quatro perfis de comportamento distintos, que determinam sua estratégia de priorização:

\begin{itemize}
    \item \textbf{Perfil Assassino:} prioriza a localização e eliminação do Wumpus, mesmo que isso implique assumir riscos maiores.
    \item \textbf{Perfil Garimpeiro}: prioriza a coleta do ouro como objetivo principal, evitando confrontos desnecessários.
    \item \textbf{Perfil Explorador}: continua explorando o ambiente mesmo após cumprir seu objetivo inicial, buscando mapear completamente o cenário.
    \item \textbf{Perfil Corajoso}: está disposto a arriscar entrar em salas consideradas perigosas, confiando em sua capacidade de reagir a imprevistos.
\end{itemize}

O \textbf{Agente Evolutivo} emprega um algoritmo genético para otimizar uma sequência de ações (cromossomo), permitindo ao usuário configurar parâmetros populacionais, taxas genéticas e uma função fitness personalizada, combinação linear de passos, coleta de ouro e sobrevivência. A cada geração, os indivíduos são executados no ambiente, e os melhores são selecionados, recombinados e mutados; ao final, o melhor cromossomo define a política de ação do agente.

\subsection{Armazenamento e Backend}

O backend da plataforma Wumpus Verse é responsável por expor a API REST e gerenciar a persistência dos dados, assegurando que todas as informações geradas pelos usuários sejam armazenadas de forma segura e consistente. Para isso, sua arquitetura integra três aspectos fundamentais: a estrutura da API, a integração direta com os agentes e os mecanismos de persistência.

Desenvolvida em Python com o framework FastAPI, a API disponibiliza endpoints que cobrem todas as funcionalidades da plataforma. Os principais endpoints estão listados na Tabela~\ref{tab:endpoints-api}, que apresenta as rotas, métodos HTTP e descrições das operações disponíveis. A comunicação entre cliente e servidor utiliza JSON como formato de dados, e a maioria das requisições é autenticada via token JWT integrado ao Firebase Authentication, garantindo que apenas usuários autorizados possam acessar e modificar seus dados.

\begin{table}[ht]
    \centering
    \small
    \begin{tabular}{|p{5.5cm}|p{2cm}|p{6.5cm}|}
        \hline
        \underline{\textbf{Rota}} & \underline{\textbf{Método}} & \underline{\textbf{Descrição}} \\
        \hline
        \texttt{/environment/list-user}, \texttt{/environment/user} & GET, POST & Lista ou cria ambientes do usuário \\
        \hline
        \texttt{/environment/execution} & POST & Executa um agente em um ambiente (recebe IDs) \\
        \hline
        \texttt{/auth/register}, \texttt{/auth/login} & POST & Registro/login, retorna token JWT \\
        \hline
        \texttt{/agents/user}, \texttt{/agents/agents} & POST, GET & Cria ou lista agentes personalizado \\
        \hline
        \texttt{/execution/user}, \texttt{/execution/list-user} & POST, GET & Salva ou lista execuções \\
        \hline
        \texttt{/execution/statics} & POST & Analisa agentes em um ambiente (estatísticas) \\
        \hline
    \end{tabular}
    \caption{Endpoints da API para gerenciamento de ambientes, agentes e execuções.}
    \label{tab:endpoints-api}
\end{table}

Complementando a estrutura da API, a integração com os agentes foi projetada para otimizar o desempenho e a confiabilidade do sistema. Como os agentes também são implementados em Python, o backend os importa diretamente como módulos, eliminando a sobrecarga e a complexidade da comunicação entre linguagens diferentes. A simulação ocorre em um ambiente controlado (classe \texttt{Environment}), que recebe o agente e o grid, executa as ações passo a passo e coleta métricas como número de passos, ouro coletado, Wumpus morto e sobrevivência. Ao final da simulação, o sistema calcula a pontuação final do agente de acordo com as regras estabelecidas. Esse desenho permite a reciclagem de código entre a API e os agentes, reduzindo significativamente a complexidade e os possíveis pontos de falha.

Por fim, no que tange à persistência de dados, a plataforma adota dois sistemas de armazenamento complementares para gerenciar diferentes aspectos das informações. O \textbf{Firebase Authentication} é responsável pelo gerenciamento de cadastro, login e redefinição de senhas de usuários, fornecendo uma infraestrutura pronta e segura que elimina a necessidade de implementar esses mecanismos do zero. O \textbf{PostgreSQL}, por sua vez, é utilizado para armazenar os dados específicos da aplicação, incluindo ambientes (grid, posições), agentes (parâmetros, tipo, função fitness) e execuções (histórico, ouros coletados, Wumpus mortos, pontuação final). A escolha do PostgreSQL garante consistência transacional e integridade referencial, características essenciais para um sistema que lida com dados complexos e inter-relacionados. Dessa forma, a combinação dessas duas ferramentas proporciona uma base robusta e escalável para todas as operações de gerenciamento de usuários e experimentos realizados na plataforma.

\subsection{Frontend}

frontend da plataforma foi desenvolvido utilizando a biblioteca JavaScript React, permitindo a construção de uma interface web dinâmica e responsiva. A arquitetura da interface foi organizada em diferentes módulos, cada um responsável por uma etapa do fluxo de utilização do sistema, desde o acesso inicial até a execução e análise dos experimentos. Esta seção descreve os principais módulos da interface, relacionando cada funcionalidade aos requisitos funcionais estabelecidos na Seção~\ref{sec:req-arq-tec}.

A \textbf{tela de apresentação} constitui a página inicial da plataforma, fornecendo acesso às principais funcionalidades disponíveis, atendendo ao requisito de disponibilidade e acesso intuitivo ao sistema. Complementarmente, a seção \textbf{Saiba Mais} reúne informações sobre o projeto, sua motivação, a equipe de desenvolvimento e uma breve contextualização sobre o problema do Mundo do Wumpus.

O sistema possui um módulo de \textbf{autenticação} composto pelas telas de login e registro, responsáveis pelo gerenciamento do acesso dos usuários. O cadastro inclui verificação de e-mail, e apenas usuários autenticados podem acessar as funcionalidades de gerenciamento e execução de experimentos. Este módulo atende indiretamente aos requisitos de persistência (RF04) ao garantir que os dados sejam associados a usuários específicos.

Para a criação e administração dos cenários, a plataforma disponibiliza as telas \textbf{Mundos} e \textbf{Editor de Mundos}, atendendo ao requisito RF01 (Ambientes Personalizados). A primeira apresenta a lista de ambientes previamente criados pelo usuário, permitindo sua visualização, edição, exclusão, importação e exportação. O editor, por sua vez, possibilita a construção de mundos bidimensionais personalizados, nos quais podem ser posicionadas as entidades do jogo, como Wumpus, poços e ouro. Embora os cenários possam assumir diferentes formatos, sua estrutura é baseada em uma malha de salas quadradas interconectadas. A Figura~\ref{fig:telas-exemplos} apresenta exemplos dessas interfaces.

\begin{figure}[ht]
    \centering
    \includegraphics[width=0.6\textwidth]{telas-plataforma.png}
    \caption{Exemplo de telas da plataforma.}
    \label{fig:telas-exemplos}
\end{figure}

A tela \textbf{Agentes} permite criar e editar agentes lógicos e evolutivos (RF02). Para agentes lógicos, o usuário seleciona perfis de comportamento (assassino, garimpeiro, explorador, corajoso); para evolutivos, define parâmetros genéticos e função fitness personalizada, podendo experimentar combinações e observar os efeitos no comportamento.

As telas \textbf{Histórico} e \textbf{Estatísticas} (RF04 e RF05) gerenciam os resultados experimentais, permitindo recuperar execuções salvas e comparar agentes por métricas como pontuação média, passos e disparos. Para iniciar uma simulação, a tela \textbf{Nova Partida} (RF03) define o mundo, o agente e a posição inicial; na tela \textbf{Execução}, o trajeto do agente é reproduzido graficamente a partir do roteiro gerado pelo backend, com controles para navegação, repetição, ajuste de parâmetros e visualização das ações (Figura~\ref{fig:tela-execucao}).

\begin{figure}[ht]
    \centering
    \includegraphics[width=0.6\textwidth]{tela-execução.png}
    \caption{Interface de execução da plataforma durante uma simulação.}
    \label{fig:tela-execucao}
\end{figure}

Além disso, o usuário pode armazenar a execução em sua conta para futuras análises, bem como exportar ou importar partidas por meio de arquivos, atendendo ao requisito RF06 (Exportação). A funcionalidade de exportação permite que os dados sejam salvos em formato JSON, enquanto a importação possibilita a reconstrução de ambientes e execuções previamente exportados.

Conforme mencionado anteriormente, a plataforma foi desenvolvida seguindo princípios de responsividade, permitindo seu acesso também por meio de dispositivos móveis. Entretanto, determinadas funcionalidades, como a criação e edição de mundos, permanecem disponíveis apenas em ambientes com teclado e mouse, uma vez que a definição da estrutura do cenário depende de interações de precisão baseadas no uso simultâneo do cursor e de comandos do teclado. Dessa forma, a adaptação para dispositivos móveis prioriza a consulta e a execução de experimentos, preservando a usabilidade das funcionalidades que demandam maior nível de interação. A Figura~\ref{fig:tela-mobile} apresenta a interface da plataforma em um dispositivo móvel, evidenciando a adaptação do layout para telas de menor dimensão.

\begin{figure}[ht]
    \centering
    \includegraphics[width=0.5\textwidth]{telas-mobile.png}
    \caption{Interface responsiva da plataforma em dispositivo
móvel.}
    \label{fig:tela-mobile}
\end{figure}